\section{Fase 4}
\setlength{\columnseprule}{0pt}
\onecolumn
\vspace{-0.3cm}
\subsection{Integración del PVA en un algoritmo de consenso}
\begin{flushleft}
\resizebox{\textwidth}{!}{%
\begin{ganttchart}[
    x unit=0.55cm, y unit title=0.65cm, y unit chart=0.65cm,
    vgrid={*1{dotted}}, hgrid,
    title height=1, title label font=\bfseries\tiny,
    bar/.style={fill=colorH},
    bar incomplete/.style={fill=colorH!25},
    group/.append style={fill=gray!50, draw=gray!50},
    progress label text={},
    bar height=0.7,
    bar top shift=0.15,
    group right shift=0.15, group top shift=0.15, group height=.7,
    milestone/.style={fill=red},
    bar label font=\small,
    group label font=\small,
    bar label node/.append style={anchor=east, align=left, text width=8.5cm},
    group label node/.append style={anchor=east, align=left, text width=8.5cm},
    bar inline label node/.append style={font=\tiny\bfseries\color{white}}]{1}{19}
  \gantttitle{Enero 2026}{4}\gantttitle{Febrero 2026}{4}\gantttitle{Marzo 2026}{4}\gantttitle{Abril 2026}{4}\gantttitle{Mayo 2026}{3}\\
  \gantttitle{S1}{1}\gantttitle{S2}{1}\gantttitle{S3}{1}\gantttitle{S4}{1}
  \gantttitle{S1}{1}\gantttitle{S2}{1}\gantttitle{S3}{1}\gantttitle{S4}{1}
  \gantttitle{S1}{1}\gantttitle{S2}{1}\gantttitle{S3}{1}\gantttitle{S4}{1}
  \gantttitle{S1}{1}\gantttitle{S2}{1}\gantttitle{S3}{1}\gantttitle{S4}{1}
  \gantttitle{S1}{1}\gantttitle{S2}{1}\gantttitle{S3}{1}\\
  \ganttgroup[inline=false]{\textbf{F4. Integración PVA-Consenso}}{1}{19}\\
  \ganttbar[progress=100,inline=false]{Seleccionar algoritmo de consenso base}{1}{2}
  \ganttbar[progress=100,inline=true]{\color{white}\textbf{100\%}}{1}{2}\\
  \ganttbar[progress=100,inline=false]{Definir como integrar consenso-PVA}{2}{4}
  \ganttbar[progress=100,inline=true]{\color{white}\textbf{100\%}}{2}{4}\\
  \ganttbar[progress=50,inline=false]{Especificar simplificaciones}{3}{5}
  \ganttbar[progress=50,inline=true]{\color{white}\textbf{50\%}}{3}{5}\\
  \ganttbar[progress=50,inline=false]{Unificar modelos y estados de los agentes}{4}{7}
  \ganttbar[progress=50,inline=true]{\color{white}\textbf{50\%}}{4}{7}\\
  \ganttbar[progress=50,inline=false]{Definir métricas para formación y consenso}{6}{9}
  \ganttbar[progress=50,inline=true]{\color{white}\textbf{50\%}}{6}{9}\\
  \ganttbar[progress=50,inline=false]{Realizar experimentación en entornos estáticos}{8}{13}
  \ganttbar[progress=50,inline=true]{\color{white}\textbf{50\%}}{8}{13}\\
  \ganttbar[progress=50,inline=false]{Realizar experimentación en entornos dinámicos}{12}{17}
  \ganttbar[progress=50,inline=true]{\color{white}\textbf{50\%}}{12}{17}\\
\end{ganttchart}
}
\end{flushleft}
